\begin{frame}{讲解：开源系统实践课真正要避免的错误路径}
  \Core{很多学生进入代码库时最大的问题不是看得太少，而是没有形成稳定阅读顺序。}
  \begin{itemize}
    \item 如果一开始只在仓库里到处搜索，很容易收集到零散函数而不是主路径。
    \item 课程更强调从入口、调度、状态、结果四个问题依次往下走。
    \item 这样做的目标不是少看代码，而是先建立可复述的系统骨架，再扩展细节。
  \end{itemize}
\end{frame}

\begin{frame}{讲解：为什么 benchmark 也是课程实践的核心材料}
  \Core{很多系统理解之所以断裂，就是因为代码导读只读 runtime，不读 workload 和结果脚本。}
  \begin{itemize}
    \item benchmark 决定了请求如何被生成、指标如何被统计、图表如何被讲述。
    \item 因此它实际上连接了代码实现与实验结论，是课程闭环里不可缺的桥梁。
    \item 真正的开源系统实践，不是把程序跑通，而是同时读懂“怎么跑”和“怎么证明”。
  \end{itemize}
\end{frame}

\begin{frame}{例子：进入一个开源 runtime 的正确姿势是什么（1）}
  \Core{开源系统实践课最怕学生一上来就“全仓库全局搜索”，最后只积累了碎片印象。}
\begin{block}{错误路径}
    先把目录全翻一遍，却没有确定任何主调用链和控制点。
  \end{block}
\end{frame}

\begin{frame}{例子：进入一个开源 runtime 的正确姿势是什么（2）}
  \Core{开源系统实践课最怕学生一上来就“全仓库全局搜索”，最后只积累了碎片印象。}
\begin{block}{正确路径}
    先找入口，再找调度循环，再找状态对象，最后找指标与脚本出口。
  \end{block}
\end{frame}

\begin{frame}{例子：进入一个开源 runtime 的正确姿势是什么（3）}
  \Core{开源系统实践课最怕学生一上来就“全仓库全局搜索”，最后只积累了碎片印象。}
\begin{block}{课程结论}
    代码实践的第一原则不是多看，而是先抓住最短的真实路径。
  \end{block}
\end{frame}

\begin{frame}{例子：为什么 benchmark 脚本也必须纳入源码导读（1）}
  \Core{系统知识不只在 runtime 本体里，benchmark 脚本经常在定义 workload、指标和报告方式。}
\begin{block}{runtime 负责执行}
    它告诉你请求如何被调度、状态如何被组织、结果如何被产生。
  \end{block}
\end{frame}

\begin{frame}{例子：为什么 benchmark 脚本也必须纳入源码导读（2）}
  \Core{系统知识不只在 runtime 本体里，benchmark 脚本经常在定义 workload、指标和报告方式。}
\begin{block}{benchmark 负责解释}
    它告诉你输入是什么、统计口径是什么、图表为什么长成这样。
  \end{block}
\end{frame}

\begin{frame}{例子：为什么 benchmark 脚本也必须纳入源码导读（3）}
  \Core{系统知识不只在 runtime 本体里，benchmark 脚本经常在定义 workload、指标和报告方式。}
\begin{block}{教学意义}
    不读 benchmark，就很容易误把“脚本默认值”当成“系统本身结论”。
  \end{block}
\end{frame}

\begin{frame}{代码例子：最小代码导读清单（1）}
  \Core{开源系统实践课适合给学生一份很短但很硬的导读清单。}
\begin{block}{导读顺序}
    {\ttfamily\small
    \detokenize{1. LLMEngine.generate()}\par
    \detokenize{2. Scheduler.schedule()}\par
    \detokenize{3. BlockManager.allocate()}\par
    \detokenize{4. bench.py / result writer}\par
    }
  \end{block}
\end{frame}

\begin{frame}{代码例子：最小代码导读清单（2）}
  \Core{开源系统实践课适合给学生一份很短但很硬的导读清单。}
\begin{block}{课堂应追问}
    为什么这四步足以构成一个最小系统理解闭环？
  \end{block}
\end{frame}

\begin{frame}{代码例子：最小代码导读清单（3）}
  \Core{开源系统实践课适合给学生一份很短但很硬的导读清单。}
\begin{block}{结论}
    好的课堂实践不会让学生读完所有代码，而是让他们先得到一条可复述的主路径。
  \end{block}
\end{frame}

\begin{frame}{代码例子：一次最小复现实验记录应该写什么（1）}
  \Core{“我跑通了”不构成复现；复现必须留下别人可重复的上下文。}
\begin{block}{记录模板}
    {\ttfamily\small
    \detokenize{model=..., precision=..., max_model_len=...}\par
    \detokenize{workload=..., warmup=..., num_requests=...}\par
    \detokenize{command=..., logs=..., result_table=...}\par
    }
  \end{block}
\end{frame}

\begin{frame}{代码例子：一次最小复现实验记录应该写什么（2）}
  \Core{“我跑通了”不构成复现；复现必须留下别人可重复的上下文。}
\begin{block}{课堂应追问}
    哪些字段缺失时，哪怕数字是对的，也不能算可复现？
  \end{block}
\end{frame}

\begin{frame}{代码例子：一次最小复现实验记录应该写什么（3）}
  \Core{“我跑通了”不构成复现；复现必须留下别人可重复的上下文。}
\begin{block}{课程结论}
    复现实验文档是工程记录，也是研究证据的一部分。
  \end{block}
\end{frame}

\begin{frame}{例子：为什么学生最容易在默认参数上踩坑（1）}
  \Core{很多所谓“复现失败”并不是机制没实现，而是默认值、环境与隐藏脚本分支没有被看见。}
\begin{block}{常见遗漏}
    温启动、随机种子、batch 上限、长度分布和精度模式。
  \end{block}
\end{frame}

\begin{frame}{例子：为什么学生最容易在默认参数上踩坑（2）}
  \Core{很多所谓“复现失败”并不是机制没实现，而是默认值、环境与隐藏脚本分支没有被看见。}
\begin{block}{常见误判}
    看到不同数字就直接怀疑代码逻辑，而不是先检查实验边界是否一致。
  \end{block}
\end{frame}

\begin{frame}{例子：为什么学生最容易在默认参数上踩坑（3）}
  \Core{很多所谓“复现失败”并不是机制没实现，而是默认值、环境与隐藏脚本分支没有被看见。}
\begin{block}{教学重点}
    进入真实开源系统后，最需要养成的是“先核对边界，再解释数字”的习惯。
  \end{block}
\end{frame}

\begin{frame}{课堂练习：给自己设计一条 30 分钟导读路线（1）}
  \Core{实践课的目标不是一口吃掉系统，而是能在短时间内稳定定位关键对象。}
\begin{block}{练习一}
    规定 30 分钟时间，写出你会先打开的 4 个文件和顺序。
  \end{block}
\end{frame}

\begin{frame}{课堂练习：给自己设计一条 30 分钟导读路线（2）}
  \Core{实践课的目标不是一口吃掉系统，而是能在短时间内稳定定位关键对象。}
\begin{block}{练习二}
    说明每个文件分别是为了解决“入口、控制、状态、证据”中的哪一个问题。
  \end{block}
\end{frame}

\begin{frame}{课堂练习：给自己设计一条 30 分钟导读路线（3）}
  \Core{实践课的目标不是一口吃掉系统，而是能在短时间内稳定定位关键对象。}
\begin{block}{练习三}
    用一句话回答：为什么 benchmark 脚本必须算作系统知识的一部分？
  \end{block}
\end{frame}
